\documentclass{article}
\usepackage{graphicx} % Required for inserting images
\usepackage{minted}
\usepackage[margin=2cm]{geometry}

\usepackage{hyperref}

\title{Tutorato Programmazione 2}
\author{Michele Porcellato, Giacomo Vettore}
\date{11 Marzo 2026}

\begin{document}

\maketitle

\section{Note}
In caso di dubbi, consultare la \href{https://docs.oracle.com/en/java/javase/25/docs/api/index.html}{documentazione Java}. Sarà accessibile anche durante l'esame.

\section{Esercizi Introduttivi}
\subsection{Media, mediana e moda}
Scrivere un programma che legga da tastiera una sequenza di $n$ numeri interi (con $n$ scelto dall'utente). Il programma deve calcolare e stampare:
\begin{itemize}
    \item \textbf{Media}: Valore decimale della somma degli elementi divisa per $n$.
    \item \textbf{Mediana}: Il valore centrale della sequenza ordinata (o la media dei due centrali se $n$ è pari).
    \item \textbf{Moda}: Il valore che compare più frequentemente.
\end{itemize}
Per ogni valore, creare una funzione apposita che prenda in input l'array e ritorni il risultato desiderato.
\paragraph{\textit{Suggerimenti:}}
\begin{itemize}
    \item Utilizzare \texttt{Arrays.sort()} usando \texttt{import java.util.Arrays} per l'ordinamento.
    \item Per la lettura dell'input usare la classe \texttt{Scanner}
    \begin{minted}{java}
Scanner sc = new Scanner(System.in);
int n = sc.nextInt(); 
    \end{minted}
\end{itemize}

\subsection{Stringa Palindroma}
Implementare un metodo statico \texttt{isPalindroma(String s)} che restituisca un valore booleano. Il programma deve leggere una stringa dall'utente e determinare se è palindroma, ignorando la differenza tra maiuscole e minuscole e gli spazi bianchi.
\begin{itemize}
    \item Esempio: "I topi non avevano nipoti" $\rightarrow$ \texttt{true}.
\end{itemize}
\paragraph{\textit{Suggerimento}:} Usare i metodi \texttt{replace()}, \texttt{toLowerCase()} per aiutarvi con l'esercizio.

\subsection{Bilanciamento delle Parentesi}
Scrivere un programma che riceva in input una stringa composta da un'espressione matematica (con sole parentesi tonde). Eg: \texttt{(5+3)*2}. Il programma deve verificare se la sequenza è bilanciata correttamente.
\begin{itemize}
    \item Esempio corretto: \texttt{(5-2)*(3*5)}
    \item Esempio errato: \texttt{3*((2+5)*(2}
\end{itemize}

\section{Esercizi OOP}
\subsection{Navi e Automobili}
Data la classe:
\begin{minted}{java}
public class Auto {
    public final String tipo;
    public Auto(final String tipo) {
        this.tipo = tipo;
    }
}
\end{minted}
Si crei una classe \textbf{Nave} che permetta di caricare automobili.
\begin{itemize}
    \item Questa classe deve definire un metodo caricaAuto([…]) a cui sarà passato un oggetto Auto.
    \item Ogni
oggetto auto verrà immagazzinato in un array d’istanza chiamato autoArray, nel primo posto
disponibile dell’array.
\end{itemize}

\subsection{Microscopi e Batteri}
Si implementi una classe \textbf{Microscopio} con un attributo \textbf{TipoMicroscopio} (elettronico o ottico) definito da un enum.
La classe dovra avere un costruttore di default (\texttt{tipoMicroscopio = TipoMicroscopio.OTTICO}) e un costruttore specifico.
Si implementi inoltre una classe Batterio con le seguenti caratteristiche:
\begin{itemize}
    \item Stringa identificativa \texttt{descrizione}
    \item La sua morfologia (Cocco, Bacillo o Spirillo) e la sua forma (rispettivamente sferica, a bastoncino e spirale). Si consiglia di usare un enum per questi valori
    \item Fare l'override di \texttt{toString()} per la stampa a video.
\end{itemize}
Infine si implementi il metodo \texttt{osserva(Batterio b)} nella classe Microscopio, che stampi a video le caratteristiche del batterio. 

\noindent \href{https://raw.githubusercontent.com/alterax05/Tutorati-P2/refs/heads/main/tutorato%201%2011-03/1TutP2/src/es2_OOP/Main.java}{\textit{File main.java}}

\subsubsection*{Esempio}
\begin{minted}{java}
Microscopio m = new Microscopio(TipoMicroscopio.ELETTRONICO);
Batterio b = new Batterio(Morfologia.COCCO);
m.osserva(b);
// Output
// Microscopio elettronico: Streptococco, Morfologia: Cocco, Forma: Sferica
\end{minted}



\subsection{Settings}

Si implementi una classe \textbf{Settings} utilizzando il \textit{Singleton Pattern}.
L'utente deve poter settare alcune impostazioni tramite opportuni metodi:
\begin{itemize}
    \item \texttt{setNetwork(Network n)} accetta un oggetto di tipo \texttt{Network}. \\ Esso ha due campi string \texttt{IP} e \texttt{SSID} (nome della rete WIFI). Inoltre c'è un campo status rappresentato dall'enum \texttt{Status} con varianti \texttt{(CONNECTED, DISCONNECTED, CONNECTING, ERROR)}
    \item \texttt{setTheme(Theme theme)} che prende un enum Theme con varianti \texttt{(DARK, LIGHT, SYSTEM)}
\end{itemize}
Inoltre sarà necessario implementere l'override a \texttt{toString()} per poter stampare sulla console le impostazioni attuali usando questo formato:  \\ \\
\texttt{Setting: \\
Network: <STATUS> <IP> <SSID> (stampare i dati solo se status = CONNECTED) \\
Theme: <THEME>
}

\noindent \href{https://raw.githubusercontent.com/alterax05/Tutorati-P2/refs/heads/main/tutorato%201%2011-03/1TutP2/src/es3_OOP/Main.java}{ \textit{File main.java}}

\end{document}
